\documentclass[12pt]{amsart}
\usepackage{graphicx, url, verbatim, systeme}
\usepackage{jdr-style, jdr-macros, jdr-linalg}
\usepackage[homework]{jdr-hwexam}
\usepackage{hyperref}
\usetikzlibrary{decorations.markings}

%\answersfalse
\answerstrue

\tikzset{
  mid arrow/.style={
    postaction={
      decoration={
        markings,
        mark=at position #1 with {\arrow{Stealth[scale=1]}},
      },
      decorate
    },
  },
  rmid arrow/.style={
    postaction={
      decoration={
        markings,
        mark=at position #1 with {\arrowreversed{Stealth[scale=1]}},
      },
      decorate
    }
  },
  mid arrow/.default={0.5},
  rmid arrow/.default={0.5},
}


\title{Interactive Linear Algebra, Exercises}
\name{Solutions}


\begin{document}

\subsection{8.1, The definition of an abstract vector space}
\begin{problems}

\problem Let \(P^2\) be the set of polynomials exactly of degree \(2\). Is this an abstract vector space? Verify the axioms or provide a counterexample.

\begin{solution}
    This is not an abstract vector space, as the zero polynomial is not of degree \(2\) and thus not included in the set.
\end{solution}

\problem Is the set of polynomials with \textbf{integer} coefficients of degree at most \(d\) an abstract vector space over \(\mathbb{R}\)? Verify the axioms or provide a counterexample.

\begin{solution}
    No, consider \(p(x) = x^d\). If we multiply it by \(\frac{1}{2}\), then the resulting polynomial no longer has integer coefficients, so the set is not closed under scalar multiplication. 
\end{solution}

\problem Is \(P_d\) still an abstract vector space if we replace addition with multiplication? What about \(P\)?

\problem Is the set of convergent sequences of real numbers an abstract vector space? Verify the axioms or provide a counterexample.

\problem Is the set of sequences of real numbers converging to a positive number an abstract vector space? Verify the axioms or provide a counterexample.









\end{problems}



\end{document}



























OLD from 2.2

\problem %In each case, write an augmented matrix in reduced row echelon form that corresponds 
%to the specified system of linear equations.
	\begin{subproblems}
	\subproblem Write an augmented matrix in RREF with three pivots that represents an 
	inconsistent system of three equations in four variables.
	\subproblem Is there a system of three equations in two variables that has infinitely many 
	solutions?  If 
	your answer is yes, write an augmented matrix in RREF that represents one such system. 
	\subproblem Are there three lines in the $xy$-plane (not necessarily distinct lines) whose 
	common intersection is exactly one 
	point?  If so, give an example and write the corresponding augmented matrix.
	\end{subproblems}
	%
	\begin{solution}
	\begin{subproblems}
	\subproblem Many possibilities.  For example,
	\[\amat{1 0 0 0 0; 0 1 0 0 0; 0 0 0 0 1}, \qquad \amat{0 1 0 0 0; 0 0 1 0 0; 0 0 0 0 1},
	\quad  \textup{ or } \quad \amat{0 0 1 0 0; 0 0 0 1 0; 0 0 0 0 1}.\]
	\subproblem Yes.  For example, consider the system
	\[ x + y = 1\]
	\[ 2x + 2y = 2\]
	\[ 3x + 3y = 3.\]
	Each equation represents the line $x + y = 1$, so the solution set is all points $(x,y)$ on the line 
	$y = 1-x$. The corresponding augmented matrix and its
	RREF are
	\[\amat{1 1 1; 2 2 2; 3 3 3} \xrightarrow[]{RREF} \amat{1 1 1; 0 0 0; 0 0 0}.\]
	\subproblem Yes.  For example, take the lines $x - y = 0$,
    $2x-2y = 0$, and $y =1$.  The first two lines are the same line,
    which intersects the non-parallel line $y=1$ at exactly one point,
    namely when $x=y=1$. To illustrate the example, we solve the
    augmented system below to find the unique solution.
	\begin{align*}
	\amat{1 -1 0; 2 -2 0;  0 1 1} 
	& \xrightarrow[]{R_2 = R_2 -2R_1}
	\amat{1 -1 0; 0 0 0; 0 1 1} \xrightarrow[]{R_1 \leftrightarrow R_2}
	\amat{1 -1 0; 0 1 1; 0 0 0} \xrightarrow[]{R_1 = R_1 + R_2} 
	\amat{1 0 1; 0 1 1; 0 0 0}.
	\end{align*}
	\end{subproblems}
	\end{solution}
\end{problems}








